%% ==========================================================================
%%  1  Introduction
%% ==========================================================================
\section{Introduction}
\label{sec:intro}
Discrete fair division studies the allocation of indivisible resources among agents with heterogeneous preferences, and lies at the interface of mathematical economics and computer science~\cite{BCE16,End17}. A central fairness criterion is \emph{envy-freeness}~\cite{Fol67,Var74}: no agent prefers another agent's bundle to her own. For indivisible items, however, an envy-free allocation need not exist, as the standard one-item two-agent instance illustrates. This has motivated a large literature on relaxations of envy-freeness, including envy-freeness up to one good (EF1)~\cite{Bud11,LMMS04} and the stronger notion of envy-freeness up to any good (EFX)~\cite{CKMPSW19}.

A second and complementary line of work retains exact envy-freeness by allowing monetary subsidies. In addition to her allocated bundle $A_i$, each agent $i$ receives a subsidy $p_i\geq 0$ from a third party, and the outcome is envy-free if
\[v_i(A_i)+p_i \geq v_i(A_j)+p_j
    \qquad\text{for all }i,j.\]
Halpern and Shah~\cite{HS19} show that every instance admits an envy-freeable allocation and study the minimum subsidy required to achieve envy-freeness. Under the standard normalization that each item value is at most one,\footnote{Such a normalization is necessary when stating absolute bounds on monetary subsidies.} they show that, for a given envy-freeable allocation, a total subsidy of at most $(n-1)m$ may be required in the worst case, and conjecture that a total subsidy of $n-1$ always suffices when the allocation can also be chosen. Brustle et al.~\cite{BDNSV20} prove this conjecture for additive valuations, obtaining an allocation that is also EF1 and balanced, and show that for general monotone valuations a subsidy of at most $2(n-1)$ per agent suffices.

Barman et al.~\cite{BKNS22} obtain a substantially stronger guarantee for dichotomous valuations. When every marginal $v_i(S\cup\{g\})-v_i(S)$ belongs to $\{0,1\}$, they show that there exists an allocation admitting subsidies in $\{0,1\}^n$, with total subsidy at most $n-1$, and that such an outcome can be computed in polynomial time in the value-oracle model. This bound is tight: for $n$ agents and a single good valued at one by every agent, the $n-1$ agents who do not receive the good each require one unit of subsidy.

In this paper, we consider the corresponding problem for indivisible chores, where receiving an additional item weakly decreases an agent's utility. Chore-allocation problems arise naturally in applications such as shift assignment, on-call duties, teaching and refereeing loads, committee service, and maintenance tasks. Monetary compensation is also natural in such settings, where it may be interpreted as compensation for undertaking an undesirable task or additional workload.

We study \emph{negative dichotomous} valuations: for every agent $i$, every $S\subseteq \items$, and every $g\in \items\setminus S$, \[ v_i(S\cup\set{g})-v_i(S)\in\set{0,-1}. \] Thus, every item is a chore for every agent, and adding a chore changes an agent's utility by either zero or minus one. This domain is the sign analogue of the dichotomous goods model of~\cite{BKNS22}.

Binary marginals capture settings in which an additional task either creates a new unit of burden or can be absorbed by an existing commitment. For example, an agent already covering an on-call window may incur no additional cost from another incident during that window, or a driver already routed through a district may absorb an additional stop at no extra cost. Since whether a chore is costly may depend on the rest of the assigned bundle, the model naturally allows non-additive valuations, including valuations that are not subadditive.

Although chore allocation is formally related to goods allocation by reversing the direction of preferences, this correspondence does not generally extend to algorithmic arguments. Several works have observed that algorithms and structural results for goods need not carry over directly to chores~\cite{aziz2017algorithms,sun2023fairness, bhaskar2020approximate}. The same difficulty arises in our setting. The algorithm of Barman et al.~\cite{BKNS22} for dichotomous goods is incremental: it allocates the goods one at a time while maintaining an envy-freeable partial allocation whose required subsidies lie in $\{0,1\}^n$. For chores, one might naturally try to preserve the same invariant by assigning each new chore to an agent who currently receives no subsidy. This approach, however, may not succeed. There are negative dichotomous instances in which a partial allocation admits an envy-free solution with subsidies in $\{0,1\}^n$, but after introducing one additional chore, no assignment of that chore to an agent preserves the unit-subsidy guarantee; this remains true even if the previously allocated bundles are allowed to be permuted before assigning the new chore. We give such an example in Appendix~\ref{app:extension-failure}. Thus, the one-item extension property underlying the goods-side construction fails for chores, and the argument of Barman et al.~\cite{BKNS22} does not directly carry over.


A general subsidy bound for our setting follows from the work of Kawase et al.~\cite{kawase2025towards}. They show that, for doubly monotone valuations with marginal values bounded in absolute value by one, an EF1 allocation can be converted into an envy-free solution with a subsidy of at most $n-1$ per agent and at most $n(n-1)/2$ in total. Negative dichotomous valuations form a subclass of this domain. Thus, Kawase et al.~\cite{kawase2025towards} already imply, for our setting, a subsidy bound that is independent of the number of chores: at most $n-1$ per agent and $n(n-1)/2$ in total. The question is whether this can be strengthened to the unit-subsidy guarantee known for dichotomous goods. Our main result answers this question affirmatively.

\begin{theorem*}
For every instance with negative dichotomous valuations, there exists a complete allocation $\mathcal A$ and a subsidy vector
$p\in\{0,1\}^n$ such that $(\mathcal A,p)$ is envy-free and
\[
    \sum_{i\in N} p_i\le n-1.
\]
Moreover, $\mathcal A$ is EF1 even before subsidies are paid, and such an allocation and subsidy vector can be computed in polynomial time in the value-oracle model.
\end{theorem*}

The total subsidy bound is tight, even for identical binary additive costs. Consider $n$ agents and $n-1$ chores with $c_i(S)=|S|$ for every agent $i$. In every complete allocation, some agent receives the empty bundle. Comparing each agent with such an agent in the envy-freeness inequalities gives $p_i\ge |A_i|$, and therefore
\[
\sum_{i\in N}p_i\ge \sum_{i\in N}|A_i|=n-1.
\]


Our proof follows a different route from the incremental construction of Barman et al.~\cite{BKNS22}. We start from the polynomial-time algorithm of Tao et al.~\cite{tao2025existence}, which returns an envy-free partial allocation while leaving at most $n-1$ chores unassigned. At termination, the remaining chores are assigned to distinct agents in a suitable tail strongly connected component of the associated equality graph. We then identify the subsidised agents through a backward-closure operation in this graph. The resulting complete allocation is EF1, and assigning one unit of subsidy to the agents in this closure makes it exactly envy-free. 

The unit-subsidy guarantee cannot, in general, be strengthened to require Pareto optimality: even with two agents and identical binary submodular costs, every Pareto-optimal allocation may require a subsidy of \(2\) for some agent (see Proposition \ref{prop_1}).

Section~\ref{sec:prelim} introduces the model and the necessary preliminaries. Section~\ref{sec:main} presents the algorithm and proves the main result.


% \paragraph{The baseline and the lower bound.}
% Two observations provide useful benchmarks for the problem, and we state them explicitly so that the contribution can be measured against them. On the one hand, the existence of a bounded subsidy is not in question. Negative dichotomous valuations are doubly monotone --- every item is a chore for every
% agent, which is a special case of every item being a good or a chore for each
% agent --- and they satisfy the normalisation
% $\abs{v_i(S \cup \set{e}) - v_i(S)} \le 1$ that~\cite[\S2]{kawase2025towards} imposes. The
% reduction of Kawase et al., which converts any EF1 allocation into an envy-free
% one, therefore applies directly and yields a subsidy of at most $n-1$ per
% agent and $n(n-1)/2$ in total~\cite[Theorem~1]{kawase2025towards}. Two details concerning this application are worth noting. Their normalisation is two-sided, so
% it bounds chore disutility as well as value; and the EF1 they require is the
% two-sided form of Definition~\ref{def:ef1}, which removes one item from
% \emph{either} bundle, not the goods-only form
% of~\cite{HS19} --- the two differ precisely when chores are present. An EF1
% allocation is guaranteed to exist and to be computable in polynomial time for
% doubly monotone valuations, so the reduction has an input; for that
% step~\cite[\S2]{kawase2025towards} cites~\cite{LMMS04} and~\cite{BSV21}. In particular, while ~\cite{LMMS04} concerns goods, ~\cite{BSV21} provides the corresponding existence result required in the chore setting.

% On the other hand, the lower bound of the goods model survives intact.

% \begin{example}[The $n-1$ lower bound transfers]
% \label{ex:lowerbound}
% Take $n$ agents and $m = n-1$ chores, with $\cost_i(S) = \abs{S}$ for every
% agent $i$ --- identical, binary, additive costs, which are in particular
% negative dichotomous. Every complete allocation leaves at least one agent
% empty-handed. Under the balanced allocation, in which $n-1$ agents hold one
% chore each and one agent holds nothing, each of the $n-1$ loaded agents envies
% the empty agent by exactly $1$ and no other envy is present, so the minimum
% subsidy assigns $1$ to each loaded agent and $0$ to the empty one, for a total of
% $n-1$. No allocation does better: consider any complete allocation $\mathcal{A}=(A_1,\ldots,A_n)$, and write $s_i=|A_i|$. Since there are only $n-1$ chores, $\sum_{i=1}^n s_i=n-1$, and therefore at least one agent, say $h$, receives no chore: $s_h=0$. Suppose that a subsidy vector $p\in\mathbb{R}_+^n$ makes the allocation envy-free. In cost form, envy-freeness requires
% \[
% c_i(A_i)-p_i
% \le
% c_i(A_j)-p_j
% \qquad
% \text{for all } i,j.
% \]
% Comparing every agent $i$ with the empty-bundle agent $h$, we obtain
% \[s_i-p_i=c_i(A_i)-p_i\le c_i(A_h)-p_h=-p_h.\]
% Hence
% \[
% p_i\ge s_i+p_h\ge s_i
% \qquad
% \text{for every } i.
% \]
% Summing over all agents gives
% \[
% \sum_{i=1}^n p_i
% \ge
% \sum_{i=1}^n s_i
% =
% n-1.
% \]
% \end{example}

% So the target statement is precisely the analogue of~\cite[Theorem~4]{BKNS22}.

% \ifworking
% \begin{conjecture}
% \label{conj:main}
% For every instance with negative dichotomous valuations there is an envy-free
% solution $(\alloc, \subsidy)$ with $\subsidy \in \set{0,1}^n$, hence with total
% subsidy at most $n-1$; the allocation $\alloc$ is moreover EF1 on its own,
% before any subsidy is paid; and both can be computed in polynomial time given
% value-oracle access.
% \end{conjecture}
% \else
% \begin{theorem}[Main result]
% \label{conj:main}
% For every instance with negative dichotomous valuations, there is an envy-free
% solution $(\alloc, \subsidy)$ with $\subsidy \in \set{0,1}^n$, hence with total
% subsidy at most $n-1$; the allocation $\alloc$ is moreover EF1 on its own,
% before any subsidy is paid; and such an allocation and subsidy vector can be computed in polynomial time given value-oracle access.
% \end{theorem}
% \fi

% Example~\ref{ex:lowerbound} shows this bound is tight. The guarantee of \cite{kawase2025towards} is $n-1$ per agent and $n(n-1)/2$ in total, whereas our target is at most one per agent and $n-1$ in total, yielding a linear-factor improvement.


% \paragraph{Our contribution.}
% We establish the statement above, for every number of agents\ifworking\
% (\Cref{sec:approach13})\else\ (Theorem~\ref{thm:m-main})\fi. The allocation
% produced is EF1 in the sense of Definition~\ref{def:ef1chores} before any money
% changes hands, and the addition of subsidies upgrades the resulting allocation to exact envy-freeness. The
% proof takes as
% input an envy-free \emph{partial} allocation --- the algorithm of Tao et
% al.~\cite{tao2025existence} produces one in polynomial time, leaving at most $n-1$ chores
% unassigned --- and completes it: the residual chores are assigned to distinct
% agents inside the strongly connected component at which that algorithm halts,
% and the paid set is the backward closure of those recipients under the equality
% relation $\cost_i(X_i) = \cost_i(X_j)$. Envy-freeness of the completion is
% verified directly, so no appeal to the characterisation of envy-freeable
% allocations is required. \Cref{sec:prelim} fixes the model and records what
% transfers from the goods theory unchanged\ifworking.\else; \Cref{sec:main}
% contains the proof.\fi

\subsection{Related work}

The literature on envy-freeness with subsidies has been developed primarily for goods. Halpern and Shah~\cite{HS19} establish the basic characterisation of envy-freeable allocations and obtain the optimal total subsidy bound of $n-1$ for binary additive valuations. Goko et al.~\cite{goko2024fair} extend the unit-subsidy guarantee to binary submodular valuations, together with strategy-proofness, while Barman et al.~\cite{BKNS22} show that no structure beyond binary marginals is required. The result most directly relevant to our setting is due to Kawase et al.~\cite{kawase2025towards}, who consider doubly monotone valuations and show that any EF1 allocation can be converted into an envy-free solution using at most $n-1$ subsidy per agent and $n(n-1)/2$ in total. Since negative dichotomous valuations form a subclass of doubly monotone valuations, this provides the previously known $m$-independent benchmark for our problem.

Fair allocation of indivisible chores without subsidies has also received considerable attention. Bhaskar et al.~\cite{BSV21} give polynomial-time algorithms for EF1 allocations of chores and doubly monotone instances, and show that deciding the existence of an exactly envy-free chore allocation is NP-complete already for binary additive costs. Tao et al.~\cite{tao2025existence} study binary-marginal chores and prove the existence of an EFX and Pareto-optimal allocation for binary additive costs. The additive assumption cannot be dropped in general: Lin et al.~\cite{LLTZ26} construct binary XOS and binary supermodular instances for which no complete EFX allocation exists. Our result is different in that it guarantees exact envy-freeness after subsidies throughout the full binary-marginal chore domain, while the underlying allocation is always EF1.

Subsidies have also been studied directly for chore allocation. For proportionality under additive costs, tight bounds are known for equal entitlements~\cite{wu2025revisiting}, and weighted variants have been studied in~\cite{wu2024tree,GSW25}. More recently, Lu et al.~\cite{LMS26} show that for additive instances containing both goods and chores, one unit of subsidy per agent suffices and the total bound of $n-1$ is tight. Their result, however, does not subsume the present setting: negative dichotomous valuations need not be additive. Conversely, additive valuations need not be dichotomous. Thus, prior to the present work, the unit-subsidy guarantee was known for additive chores~\cite{LMS26} and for dichotomous goods~\cite{BKNS22}, but not for non-additive dichotomous chores.

To the best of our knowledge, this is the first unit-subsidy guarantee for envy-freeness under negative dichotomous valuations.


% \subsection{Related work}
% The subsidy line for goods is described above. Within it, the dichotomous subclass has attracted separate attention: Halpern and Shah~\cite{HS19} settle
% binary additive valuations with a total subsidy of $n-1$, Goko et
% al.~\cite{goko2024fair} obtain the analogous bound for binary submodular valuations
% together with a strategy-proof mechanism, and Barman et al.~\cite{BKNS22} subsume
% both by dropping every structural assumption beyond binary marginals. Dichotomous
% preferences are themselves a well-studied restriction in social
% choice~\cite{BMS05,BEF21}. On the computational side, Caragiannis and
% Ioannidis~\cite{CI21} show that minimising the subsidy is NP-hard and give a
% fully polynomial-time approximation scheme for a constant number of agents, and
% Narayan et al.~\cite{NSV21} study the welfare cost of achieving envy-freeness
% with transfers.

% The work of Kawase et al.~\cite{kawase2025towards} is particularly relevant to the present setting because it handles general doubly monotone valuations, and hence includes non-additive chore valuations. Working with doubly monotone valuations under the two-sided normalisation,
% it separates the allocation and subsidy steps that are handled jointly in~\cite{BDNSV20}: rather than constructing a specific allocation and showing that it is cheap to subsidise, it takes \emph{any} EF1
% allocation as input and bounds the subsidy needed to convert it, obtaining $n-1$ per agent and $n(n-1)/2$ in total. Since EF1 allocations are known to exist for doubly monotone valuations, this yields the stated subsidy guarantee throughout that class. This result, therefore, provides the natural benchmark against which we compare our subsidy bound.

% A complementary approach dispenses with subsidies and instead considers approximate notions of envy-freeness.  Bu et al.~\cite{BSY23} show that EFX allocations always
% exist and can be computed in polynomial time for dichotomous goods valuations, thus
% extending the matroid-rank result of Babaioff et al.~\cite{BEF21}. These results provide two complementary guarantees on the goods side: either pay $0$ or $1$ per agent and obtain exact envy-freeness~\cite{BKNS22}, or pay nothing and obtain EFX~\cite{BSY23}. This paper establishes the first branch for
% chores.

% \paragraph{The chores side.}
% Subsidies have also been studied in the allocation of indivisible chores, although much of this literature concerns fairness notions
% other than envy-freeness. Subsidies have, in particular, been studied for achieving proportionality under additive costs. For agents with equal entitlements, tight subsidy bounds are known: the optimal total subsidy is $n/4$ when $n$ is even and $(n^2-1)/(4n)$ when $n$ is odd~\cite{wu2025revisiting}. In the more general
% weighted setting, a total subsidy of $n/3-1/6$ is known to suffice~\cite{wu2024tree}, and Garg et al.~\cite{GSW25} obtain this guarantee together with Pareto optimality. These results concern proportionality rather than envy-freeness, and therefore do not directly bound the subsidy required in our setting.

% Turning to approximate fairness without subsidies, there is a direct chore analogue of~\cite{BSY23}. Tao et al.~\cite{tao2025existence} study chores with binary marginals and show that for binary \emph{additive} costs, an allocation that is simultaneously EFX and Pareto optimal always exists and is computable in polynomial time. So on the chores side the second half of the binary dichotomy --- pay nothing, obtain EFX --- is
% established at least for additive costs, while the first half pays $0$ or $1$ per agent and obtain exact envy-freeness, is what this paper establishes. The restriction to additive costs there cannot simply be dropped: Lin et al.~\cite{LLTZ26} exhibit instances with $18$ agents and $53$ chores, one with binary XOS costs and one with binary supermodular costs, in which no complete EFX allocation exists at all; the binary submodular case remains open. Thus, complete EFX allocations need not exist throughout the binary-marginal chore domain, whereas the allocation constructed below is always EF1, and becomes exactly envy-free once at most one unit of subsidy per agent is paid. Exact envy-freeness without subsidies is not always attainable; indeed, Bhaskar et al.~\cite{BSV21} show that deciding whether an exactly envy-free allocation of chores exists is NP-complete already for binary additive costs, and give a polynomial-time EF1 algorithm for chores and for doubly
% monotone instances.

% Recent work has also considered subsidies for additive instances involving both goods and chores. Lu et
% al.~\cite{LMS26} prove the goods-and-chores analogue of~\cite{BDNSV20}: for
% additive utilities in which each item is a good for some agents and a chore for
% others, with every $\abs{v_i(g)} \le 1$, a subsidy of one dollar per agent
% suffices, the bound is tight, and the allocation is computable in polynomial
% time --- hence $n-1$ in total. In particular, this establishes the optimal subsidy bound for additive chore instances.

% This result does not subsume the present setting, for the same reason that~\cite{BKNS22} is not a corollary of~\cite{BDNSV20}: dichotomous valuations need not be additive, and additive valuations need not be dichotomous. The relationship among these results can be summarised as follows:

% \begin{center}
% \begin{tabular}{@{}lll@{}}
% \toprule
%  & goods & chores \\
% \midrule
% additive     & \cite{BDNSV20} & \cite{LMS26} \\
% dichotomous  & \cite{BKNS22}  & \textbf{this paper} \\
% \bottomrule
% \end{tabular}
% \end{center}

% \noindent
% each of the four yielding    
  a per-agent subsidy bound of one and a total subsidy bound of $n-1$. Note also that~\cite{LMS26} does not subsume~\cite{kawase2025towards} for the valuation class considered here:
% for valuations that are dichotomous but not additive, the $n-1$ per agent
% of~\cite{kawase2025towards} was the best previously known bound, and the present result improves this by a linear factor.

% To the best of our knowledge, envy-freeness with subsidies has not previously been studied for chores under dichotomous costs.